% ECONOMICS_PROBLEM: {"id": "C62-182", "title": "Realistic equilibrium adjustment", "jel_code": "C62", "source_id": "OP-0271"}
% !TeX program = xelatex
\documentclass[11pt,a4paper]{article}
\usepackage[margin=23mm]{geometry}
\usepackage{fontspec}
\setmainfont{Latin Modern Roman}
\usepackage{amsmath,amssymb,mathtools}
\usepackage{xcolor,enumitem,booktabs,longtable,array,xurl,hyperref}
\definecolor{muted}{HTML}{53636C}
\hypersetup{colorlinks=true,urlcolor=blue,linkcolor=blue}
\setlength{\parindent}{0pt}
\setlength{\parskip}{5pt}
\newcommand{\R}{\mathbb R}
\newcommand{\E}{\mathbb E}
\newcommand{\N}{\mathbb N}
\newcommand{\ind}{\mathbf 1}
\newcommand{\Var}{\operatorname{Var}}
\newcommand{\Cov}{\operatorname{Cov}}
\newcommand{\supp}{\operatorname{supp}}
\DeclareMathOperator*{\argmin}{arg\,min}
\DeclareMathOperator*{\argmax}{arg\,max}
\newcommand{\entrymeta}[3]{{\sffamily\small\color{muted}\textbf{\texttt{#1}}\quad|\quad #2\par #3\par}}
\newcommand{\Problemref}[1]{\texttt{#1}}
\newcommand{\JELref}[2]{\texttt{#2} (JEL \texttt{#1})}
\begin{document}
\section*{Realistic equilibrium adjustment}
\textbf{Problem ID:} \texttt{C62-182}\quad \textbf{JEL:} \texttt{C62}\par
% BEGIN ECONOMICS STATEMENT
\label{problem:OP-0271}
{\sffamily\small\color{muted}Original subject group: Economic theory, equilibrium and social choice\par}
\label{definitions:problem:30007552}
\textbf{Audit classification:} Background; proposed extension. The cited paper proves stability results for a different forward-looking model. The delayed rule is an editorial extension; no published statement or residual unsolved status was established.
\subsection*{Definitions and precise research target}
\textbf{Model and research target.} Consider a finite pure-exchange economy with agents \(i=1,\ldots,n\), goods \(g=1,\ldots,m\), positive endowments \(\omega_i\), and strictly increasing, strictly concave, twice continuously differentiable utilities \(u_i\). Assume demands are unique at every strictly positive normalized price. Let \(\Delta_{++}\) be the positive unit price simplex, \(z(p)\) aggregate excess demand, and \(p^*\in\Delta_{++}\) a competitive equilibrium satisfying \(z(p^*)=0\). Fix a delay \(\tau\geq0\), diagonal positive adjustment matrix \(D\), and projection \(\Pi\) onto the simplex tangent space. Study the proposed decentralized rule
\[
 \dot p(t)=\Pi D z(p(t-\tau)),\qquad p(s)=\phi(s)\quad(-\tau\leq s\leq0),
\]
where \(\phi\) is a continuous positive initial price history. Restrict attention to economies and histories for which the resulting trajectory remains in a specified compact subset of \(\Delta_{++}\). Determine economically interpretable necessary or sufficient restrictions on preferences, endowments, \(D\), and \(\tau\) guaranteeing convergence to equilibrium for every admissible history. Derive sharp delay bounds where possible and construct economies demonstrating their failure outside the guaranteed region. Both local stability and genuinely global convergence must be distinguished.

\emph{Definitions and maintained assumptions (editorial unless a source definition is named).} Let \(\mathbb R_+^m\) be the consumption set, \(\Delta_{++}=\{p\in\mathbb R^m:p_g>0,\ \sum_g p_g=1\}\), and
\[x_i(p)=\arg\max\{u_i(x):x\in\mathbb R_+^m,\ p\cdot x\le p\cdot\omega_i\},\qquad z(p)=\sum_i(x_i(p)-\omega_i).\]
Assume single-valued demands and that \(z\) is locally Lipschitz on \(\Delta_{++}\); smooth utilities alone do not automatically imply a differentiable demand at consumption corners. Use the Euclidean tangent projection \(\Pi=I-\mathbf1\mathbf1^\top/m\), rather than an unspecified price-dependent projection. A history belongs to \(C([-\tau,0],\Delta_{++})\), with the supremum norm. An admissible history is one whose unique delay-equation solution exists for all \(t\ge0\) and stays in a fixed compact \(K\subset\Delta_{++}\). Global convergence means that each admissible solution has a limit \(p^\infty\) with \(z(p^\infty)=0\); it does not require the same limit for different histories. Local asymptotic stability adds Lyapunov stability in the history norm. An interior rest point of the projected equation is an equilibrium: \(\Pi Dz=0\) implies \(Dz=c\mathbf1\), and Walras' law gives \(c\sum_g p_g/D_{gg}=0\).
\subsection*{Variants and assumption changes}
\begin{enumerate}
\item \textbf{Local delays (editorial).} Assume \(z\) is \(C^1\) near an isolated equilibrium. On the tangent space analyze the roots of \(\det(\lambda I-\Pi Dz'(p^*)e^{-\lambda\tau})=0\). Determine the first delay at which a characteristic root reaches the imaginary axis; a linearization conclusion requires the usual noncritical-root hypotheses.
\item \textbf{Global domain (editorial).} Impose a stated gross-substitutes or monotonicity restriction and first prove that \(K\) is forward invariant for all histories taking values in \(K\). Then seek a uniform delay bound guaranteeing convergence on that domain. Compactness must be proved or imposed; it is not provided by tangent projection.
\end{enumerate}
\subsection*{References and source audit}
\begin{enumerate}
\item Abstract; introduction; §§5–7. Supports the stated research area; exact displayed specification is editorial. \url{https://economics.mit.edu/sites/default/files/inline-files/taton_NBER_new.pdf}
\end{enumerate}
\begin{enumerate}\raggedright
\item\hypertarget{definitions:source:30007552:1}{} Iván Werning; Guido Lorenzoni. 2026. \emph{Tâtonnement and Price Setting in General Equilibrium}. Abstract; §§5–7, especially §7.2 Proposition 7.2.
\url{https://economics.mit.edu/sites/default/files/inline-files/taton_NBER_new.pdf}.
DOI: \url{https://doi.org/10.3386/w35205}.
\end{enumerate}

\subsection*{Common conventions: Source mathematical conventions}

These conventions apply to each entry unless explicitly replaced.

\textbf{Models and identification.} A primitive model \(m\) specifies all probability laws, feasible choices, information, equilibrium and measurement rules used by its target. The admissible class \(\mathcal M\) is fixed independently of outcomes. If \(P_O(m)\) is its observable probability law and \(\tau(m)\) the target, its identified set at observable law \(P\) is
\[
 \mathcal I_\tau(P)=\{\tau(m):m\in\mathcal M,\ P_O(m)=P\}.
\]
Point identification means that this set is a singleton. An empty set signals incompatibility of data and assumptions. Coordinate bounds need not describe the joint set. Bounds are sharp only if every retained target is attainable or arbitrarily approachable by compatible models. Finite bounds require explicit restrictions on unobserved quantities; unrestricted confounding, hidden wealth, extrapolation or arbitrary equilibrium selection can yield uninformative sets.

\textbf{Causal interventions.} A potential outcome \(Y_i(p)\) is the outcome under a complete policy regime \(p\), including timing, financing, information and market spillovers. The target population and units are fixed before treatment. With interference, use the complete assignment vector or a justified exposure mapping; individual no-interference notation alone cannot identify an economy-wide intervention. A difference of mean potential outcomes does not require the cross-world joint distribution. Conditioning on post-treatment migration, survival, enrollment or employment changes the estimand and needs separate justification.

\textbf{Statistical statements.} An estimator, confidence set, forecast or test specifies a sampling law, model class and loss/coverage criterion. Uniform \((1-\alpha)\)-coverage means \(\inf_{m\in\mathcal M}\Pr_m\{\tau(m)\in C_n\}\geq1-\alpha\), or an explicitly asymptotic version. The whole parameter space trivially has coverage, so informativeness requires a stated width, risk or power objective. A forecasting improvement is relative to a named benchmark, information set and evaluation population.

\textbf{Equilibrium and optimization.} An equilibrium comprises feasible plans, optimality, beliefs where needed, budget constraints and all market/resource clearing. The primitives or a stated benchmark define each correspondence \(\mathcal E(p;m)\). Nonemptiness is assumed or proved, never inferred just from compact policy sets. A selected-equilibrium target uses a declared selection rule; a robust target quantifies over all permitted equilibria. Use suprema or infima unless attainment is established. An \(\varepsilon\)-optimal policy is feasible and within \(\varepsilon\) of the finite optimum value. Report an empty feasible set separately.

\textbf{Welfare and accounting.} Normative utility, interpersonal normalization, social weights, numeraire, discounting and the use of public receipts are primitives. Behavioral utility need not coincide with the welfare criterion. Household welfare includes ownership income and rents once; adding profits again to compensating variation generally double-counts them. A monetary equivalent is defined by an explicit transfer and price convention, with monotonicity and range conditions. Average effects, mechanism contributions, transfers and welfare losses are different targets. Infinite sums require convergence and dynamic feasible plans require terminal or no-Ponzi conditions.

\textbf{Source versions.} A working-paper date, revision date and journal publication year are distinguished. A DOI for a journal version is not automatically the DOI of an earlier manuscript. Locators refer to the stated version and printed pages where specified. A metadata-only check does not verify a claim about a full-text conclusion. Every entry's source subsection narrows the provenance claim accordingly.
% END ECONOMICS STATEMENT
\end{document}
